\section{Index-keyed families}

\import{set/choice.tex}
\import{set/cons.tex}

An indexed family retains an explicit index set, so equal fibres at distinct
indices remain distinct coordinates.  Its dependent product consists of the
functions that choose one element from each indexed fibre.

\begin{definition}\label{pairwise_disjoint_on}
    $\pairwiseDisjointOn{P}{I}$ iff for all $i,j\in I$ such that
        $i\neq j$ we have $P(i)$ is disjoint from $P(j)$.
\end{definition}

\begin{abbreviation}\label{pairwise_disjoint_on_phrase}
    $P$ is pairwise disjoint on $I$ iff $\pairwiseDisjointOn{P}{I}$.
\end{abbreviation}

\begin{definition}\label{family_remainder}
    $\familyRemainder{P}{S}=S\setminus\unions{\ran{P}}$.
\end{definition}

\begin{definition}\label{complete_disjoint_family_at}
    $\completeDisjointFamilyAt{P}{S}{I}{i_0}
        =\cons{(i_0,P(i_0)\union\familyRemainder{P}{S})}
            {\restrl{P}{I\setminus\{i_0\}}}$.
\end{definition}

\begin{proposition}\label{pairwise_disjoint_on_reindex_injection}
    Suppose $P$ is a function on $I$.
    Let $e\in\inj{J}{I}$.
    Suppose $P$ is pairwise disjoint on $I$.
    Then $P\circ e$ is pairwise disjoint on $J$.
\end{proposition}
\begin{proof}
    $e\in\funs{J}{I}$ by \obv{inj}.
    Hence $e$ is a function and $\dom{e}=J$ and $\ran{e}\subseteq I$
        by \obv{funs_is_function,funs,funs_ran}.
    $P$ is a function and $\dom{P}=I$.
    Show that for all $i,j\in J$ such that $i\neq j$
        we have $(P\circ e)(i)$ is disjoint from $(P\circ e)(j)$.
    \begin{subproof}
        Fix $i,j\in J$.
        Assume $i\neq j$.
        $e(i),e(j)\in I$ by \obv{funs_type_apply}.
        $e(i)\neq e(j)$ by \obv{inj,injective_function,funs}.
        Hence $P(e(i))$ is disjoint from $P(e(j))$
            by \obv{pairwise_disjoint_on}.
        $(P\circ e)(i)=P(e(i))$ and $(P\circ e)(j)=P(e(j))$
            by \obv{circ_apply}.
    \end{subproof}
    Follows by \obv{pairwise_disjoint_on}.
\end{proof}

\begin{proposition}\label{pairwise_disjoint_on_subdomain}
    Suppose $J\subseteq I$ and $P$ is pairwise disjoint on $I$.
    Then $P$ is pairwise disjoint on $J$.
\end{proposition}
\begin{proof}
    Follows by \obv{pairwise_disjoint_on,subseteq}.
\end{proof}

\begin{proposition}\label{family_remainder_subseteq}
    $\familyRemainder{P}{S}\subseteq S$.
\end{proposition}
\begin{proof}
    Follows by \auto{family_remainder,setminus,subseteq}.
\end{proof}

\begin{proposition}\label{family_remainder_disjoint_member}
    Suppose $P$ is a function on $I$.
    Let $i\in I$.
    Then $\familyRemainder{P}{S}$ is disjoint from $P(i)$.
\end{proposition}
\begin{proof}
    $P$ is a function and $\dom{P}=I$.
    Suppose not.
    Take $x$ such that $x\in\familyRemainder{P}{S},P(i)$
        by \auto{disjoint}.
    $P(i)\in\ran{P}$ by \auto{fun_ran_iff}.
    Hence $x\in\unions{\ran{P}}$ by \auto{unions_iff}.
    $x\notin\unions{\ran{P}}$ by \auto{family_remainder,setminus}.
    Contradiction.
\end{proof}

\begin{proposition}\label{complete_disjoint_family_at_type}
    Let $P\in\funs{I}{\pow{S}}$.
    Let $i_0\in I$.
    Then $\completeDisjointFamilyAt{P}{S}{I}{i_0}
        \in\funs{I}{\pow{S}}$.
\end{proposition}
\begin{proof}
    Let $J=I\setminus\{i_0\}$.
    $J\subseteq I$ by \auto{subseteq,setminus}.
    Hence $\restrl{P}{J}\in\funs{J}{\pow{S}}$
        by \obv{funs_restrl}.
    $i_0\notin J$ by \auto{setminus,singleton_iff}.
    $P(i_0)\in\pow{S}$ by \obv{funs_type_apply}.
    Hence $P(i_0)\subseteq S$ by \obv{powerset_elim}.
    $\familyRemainder{P}{S}\subseteq S$
        by \obv{family_remainder_subseteq}.
    Hence $P(i_0)\union\familyRemainder{P}{S}\subseteq S$
        by \auto{subseteq,union_iff}.
    Hence $P(i_0)\union\familyRemainder{P}{S}\in\pow{S}$
        by \obv{powerset_intro}.
    Hence
        $\cons{(i_0,P(i_0)\union\familyRemainder{P}{S})}
            {\restrl{P}{J}}
        \in\funs{\cons{i_0}{J}}{\pow{S}}$
        by \obv{extend_function_at_new_point}.
    $\cons{i_0}{J}=I$ by \auto{cons_remove}.
    Follows by \auto{complete_disjoint_family_at}.
\end{proof}

\begin{proposition}\label{complete_disjoint_family_at_selected}
    Let $P\in\funs{I}{\pow{S}}$.
    Let $i_0\in I$.
    Then $(\completeDisjointFamilyAt{P}{S}{I}{i_0})(i_0)
        =P(i_0)\union\familyRemainder{P}{S}$.
\end{proposition}
\begin{proof}
    Let $Q=\completeDisjointFamilyAt{P}{S}{I}{i_0}$.
    $Q$ is a function by
        \auto{complete_disjoint_family_at_type,funs_is_function}.
    $(i_0,P(i_0)\union\familyRemainder{P}{S})\in Q$
        by \auto{complete_disjoint_family_at,cons_iff}.
    Follows by \obv{function_apply_intro}.
\end{proof}

\begin{proposition}\label{complete_disjoint_family_at_other}
    Let $P\in\funs{I}{\pow{S}}$.
    Let $i_0,i\in I$.
    Suppose $i\neq i_0$.
    Then $(\completeDisjointFamilyAt{P}{S}{I}{i_0})(i)=P(i)$.
\end{proposition}
\begin{proof}
    Let $J=I\setminus\{i_0\}$.
    $P$ is a function and $\dom{P}=I$
        by \obv{funs_is_function,funs}.
    $i\in J$ by \auto{setminus,singleton_iff}.
    $(i,P(i))\in P$ by \obv{function_apply_elim}.
    Hence $(i,P(i))\in\restrl{P}{J}$ by \obv{restrl_iff}.
    Hence $(i,P(i))\in\completeDisjointFamilyAt{P}{S}{I}{i_0}$
        by \auto{complete_disjoint_family_at,cons_iff}.
    $\completeDisjointFamilyAt{P}{S}{I}{i_0}$ is a function
        by \auto{complete_disjoint_family_at_type,funs_is_function}.
    Follows by \obv{function_apply_intro}.
\end{proof}

\begin{proposition}\label{complete_disjoint_family_at_pairwise}
    Let $P\in\funs{I}{\pow{S}}$.
    Let $i_0\in I$.
    Suppose $P$ is pairwise disjoint on $I$.
    Then $\completeDisjointFamilyAt{P}{S}{I}{i_0}$
        is pairwise disjoint on $I$.
\end{proposition}
\begin{proof}
    Let $Q=\completeDisjointFamilyAt{P}{S}{I}{i_0}$.
    $P$ is a function and $\dom{P}=I$
        by \obv{funs_is_function,funs}.
    Show that for all $i,j\in I$ such that $i\neq j$
        we have $Q(i)$ is disjoint from $Q(j)$.
    \begin{subproof}
        Fix $i,j\in I$.
        Assume $i\neq j$.
        \begin{byCase}
        \caseOf{$i=i_0$.}
            $j\neq i_0$.
            $Q(i)=P(i_0)\union\familyRemainder{P}{S}$
                by \auto{complete_disjoint_family_at_selected}.
            $Q(j)=P(j)$
                by \auto{complete_disjoint_family_at_other}.
            $P(i_0)$ is disjoint from $P(j)$
                by \obv{pairwise_disjoint_on}.
            $\familyRemainder{P}{S}$ is disjoint from $P(j)$
                by \auto{family_remainder_disjoint_member}.
            Suppose not.
            Take $x$ such that $x\in Q(i),Q(j)$ by \auto{disjoint}.
            Hence $x\in P(i_0)$ or $x\in\familyRemainder{P}{S}$
                by \auto{union_iff}.
            Contradiction by \auto{disjoint}.
        \caseOf{$i\neq i_0$.}
            \begin{byCase}
            \caseOf{$j=i_0$.}
                $Q(i)=P(i)$
                    by \auto{complete_disjoint_family_at_other}.
                $Q(j)=P(i_0)\union\familyRemainder{P}{S}$
                    by \auto{complete_disjoint_family_at_selected}.
                $P(i)$ is disjoint from $P(i_0)$
                    by \obv{pairwise_disjoint_on}.
                $\familyRemainder{P}{S}$ is disjoint from $P(i)$
                    by \auto{family_remainder_disjoint_member}.
                Suppose not.
                Take $x$ such that $x\in Q(i),Q(j)$ by \auto{disjoint}.
                Hence $x\in P(i_0)$ or
                    $x\in\familyRemainder{P}{S}$ by \auto{union_iff}.
                Contradiction by \auto{disjoint}.
            \caseOf{$j\neq i_0$.}
                $Q(i)=P(i)$ and $Q(j)=P(j)$
                    by \auto{complete_disjoint_family_at_other}.
                Follows by \auto{pairwise_disjoint_on}.
            \end{byCase}
        \end{byCase}
    \end{subproof}
    Follows by \auto{pairwise_disjoint_on}.
\end{proof}

\begin{proposition}\label{complete_disjoint_family_at_cover}
    Let $P\in\funs{I}{\pow{S}}$.
    Let $i_0\in I$.
    Then $\unions{\ran{\completeDisjointFamilyAt{P}{S}{I}{i_0}}}=S$.
\end{proposition}
\begin{proof}
    Let $Q=\completeDisjointFamilyAt{P}{S}{I}{i_0}$.
    $P$ is a function and $\dom{P}=I$
        by \obv{funs_is_function,funs}.
    $Q\in\funs{I}{\pow{S}}$ by
        \auto{complete_disjoint_family_at_type}.
    Hence $Q$ is a function and $\dom{Q}=I$
        by \obv{funs_is_function,funs}.
    Show that $\unions{\ran{Q}}\subseteq S$.
    \begin{subproof}
        It suffices to show that every element of $\unions{\ran{Q}}$
            is an element of $S$ by \obv{subseteq}.
        Fix $x\in\unions{\ran{Q}}$.
        Take $B\in\ran{Q}$ such that $x\in B$ by \obv{unions_iff}.
        Take $i\in\dom{Q}$ such that $Q(i)=B$ by \obv{fun_ran_iff}.
        $Q(i)\in\pow{S}$ by \obv{funs_type_apply}.
        Follows by \auto{powerset_elim,subseteq}.
    \end{subproof}
    Show that $S\subseteq\unions{\ran{Q}}$.
    \begin{subproof}
        It suffices to show that every element of $S$
            is an element of $\unions{\ran{Q}}$ by \obv{subseteq}.
        Fix $x\in S$.
        \begin{byCase}
        \caseOf{$x\in\unions{\ran{P}}$.}
            Take $B\in\ran{P}$ such that $x\in B$ by \obv{unions_iff}.
            Take $i\in\dom{P}$ such that $P(i)=B$
                by \obv{fun_ran_iff}.
            $\dom{P}=I$ by \obv{funs}.
            Hence $i\in I$.
            \begin{byCase}
            \caseOf{$i=i_0$.}
                $Q(i_0)=P(i_0)\union\familyRemainder{P}{S}$
                    by \auto{complete_disjoint_family_at_selected}.
                Hence $x\in Q(i_0)$ by \auto{union_iff}.
                $Q(i_0)\in\ran{Q}$
                    by \auto{funs,fun_ran_iff}.
                Follows by \auto{unions_iff}.
            \caseOf{$i\neq i_0$.}
                $Q(i)=P(i)$
                    by \auto{complete_disjoint_family_at_other}.
                Hence $x\in Q(i)$.
                $Q(i)\in\ran{Q}$ by \auto{funs,fun_ran_iff}.
                Follows by \auto{unions_iff}.
            \end{byCase}
        \caseOf{$x\notin\unions{\ran{P}}$.}
            Hence $x\in\familyRemainder{P}{S}$
                by \obv{family_remainder,setminus}.
            $Q(i_0)=P(i_0)\union\familyRemainder{P}{S}$
                by \auto{complete_disjoint_family_at_selected}.
            Hence $x\in Q(i_0)$ by \auto{union_iff}.
            $Q(i_0)\in\ran{Q}$ by \auto{funs,fun_ran_iff}.
            Follows by \auto{unions_iff}.
        \end{byCase}
    \end{subproof}
    Follows by \obv{subseteq_antisymmetric}.
\end{proof}

\begin{definition}\label{dependent_product}
    $\dependentProduct{I}{A}
        =\{f\in\funs{I}{\unions{\ran{A}}}\mid
            \text{for every $i\in I$ we have $f(i)\in A(i)$}\}$.
\end{definition}

\begin{definition}\label{indexed_sum}
    $\indexedSum{I}{A}
        =\{p\in I\times\unions{\ran{A}}\mid
            \snd{p}\in A(\fst{p})\}$.
\end{definition}

\begin{definition}\label{indexed_intersection}
    $\indexedIntersection{I}{A}{U}
        =\{x\in U\mid \text{for every $i\in I$ we have $x\in A(i)$}\}$.
\end{definition}

\begin{proposition}\label{indexed_intersection_iff}
    $x\in\indexedIntersection{I}{A}{U}$ iff
        $x\in U$ and for every $i\in I$ we have $x\in A(i)$.
\end{proposition}
\begin{proof}
    Follows by \obv{indexed_intersection}.
\end{proof}

\begin{proposition}\label{indexed_intersection_subseteq}
    $\indexedIntersection{I}{A}{U}\subseteq U$.
\end{proposition}
\begin{proof}
    Follows by \obv{subseteq,indexed_intersection_iff}.
\end{proof}

\begin{proposition}\label{indexed_intersection_member_subseteq}
    Suppose $i\in I$.
    Then $\indexedIntersection{I}{A}{U}\subseteq A(i)$.
\end{proposition}
\begin{proof}
    Follows by \obv{subseteq,indexed_intersection_iff}.
\end{proof}

\begin{proposition}\label{dependent_product_iff}
    $f\in\dependentProduct{I}{A}$ iff
        $f\in\funs{I}{\unions{\ran{A}}}$ and
        for every $i\in I$ we have $f(i)\in A(i)$.
\end{proposition}
\begin{proof}
    Follows by \obv{dependent_product}.
\end{proof}

\begin{theorem}\label{choice_function_iff_dependent_product}
    $f$ is a choice function for $F$ iff
        $f\in\dependentProduct{F}{\identity{F}}$.
\end{theorem}
\begin{proof}
    $\ran{\identity{F}}=F$ by \obv{id_ran}.
    Show that if $f$ is a choice function for $F$, then
        $f\in\dependentProduct{F}{\identity{F}}$.
    \begin{subproof}
        Assume $f$ is a choice function for $F$.
        Hence $f$ is a function on $F$ and for every $A\in F$ we have
            $f(A)\in A$ by \obv{choice_function}.
        Thus $f$ is a function and $F=\dom{f}$.
        For every $A\in\dom{f}$ we have
            $f(A)\in\unions{\ran{\identity{F}}}$ by
            \auto{unions_iff}.
        Hence $f$ is a function to
            $\unions{\ran{\identity{F}}}$.
        Hence $f\in\funs{F}{\unions{\ran{\identity{F}}}}$ by
            \obv{funs_intro}.
        For every $A\in F$ we have
            $f(A)\in(\identity{F})(A)$ by \obv{id_apply}.
        Follows by \obv{dependent_product_iff}.
    \end{subproof}
    Show that if $f\in\dependentProduct{F}{\identity{F}}$, then
        $f$ is a choice function for $F$.
    \begin{subproof}
        Assume $f\in\dependentProduct{F}{\identity{F}}$.
        Hence $f\in\funs{F}{\unions{\ran{\identity{F}}}}$ and
            for every $A\in F$ we have
                $f(A)\in(\identity{F})(A)$ by
            \obv{dependent_product_iff}.
        Hence $f$ is a function and $F=\dom{f}$ by
            \obv{funs_is_function,funs}.
        Hence $f$ is a function on $F$.
        For every $A\in F$ we have $f(A)\in A$ by \auto{id_apply}.
        Follows by \obv{choice_function}.
    \end{subproof}
    Follows by \auto.
\end{proof}

\begin{proposition}\label{empty_dependent_product}
    $\dependentProduct{\emptyset}{A}=\{\emptyset\}$.
\end{proposition}
\begin{proof}
    Show that every element of $\dependentProduct{\emptyset}{A}$ is an
        element of $\{\emptyset\}$.
    \begin{subproof}
        Fix $f\in\dependentProduct{\emptyset}{A}$.
        $f\in\funs{\emptyset}{\unions{\ran{A}}}$ by
            \obv{dependent_product_iff}.
        Hence $f$ is a function and $\dom{f}=\emptyset$ by
            \obv{funs_is_function,funs}.
        Show that every element of $f$ is an element of $\emptyset$.
        \begin{subproof}
            Fix $w\in f$.
            Take $x,y$ such that $w=(x,y)$ by \obv{relation}.
            Hence $x\in\dom{f}$ by \auto{dom_iff}.
            Contradiction by \obv{emptyset}.
        \end{subproof}
        Hence $f\subseteq\emptyset$ by \obv{subseteq}.
        Hence $f=\emptyset$ by \obv{subseteq_emptyset_iff}.
        Follows by \obv{singleton_iff}.
    \end{subproof}

    $\emptyset$ is a function to $\unions{\ran{A}}$ by
        \obv{codom_of_emptyset_can_be_anything}.
    $\dom{\emptyset}=\emptyset$ by \obv{dom_emptyset}.
    Hence $\emptyset\in\funs{\emptyset}{\unions{\ran{A}}}$ by
        \obv{funs_intro}.
    For every $i\in\emptyset$ we have $\apply{\emptyset}{i}\in A(i)$ by
        \auto{emptyset}.
    Hence $\emptyset\in\dependentProduct{\emptyset}{A}$ by
        \obv{dependent_product_iff}.
    Hence $\{\emptyset\}\subseteq\dependentProduct{\emptyset}{A}$ by
        \obv{singleton_subset_intro}.
    Follows by \obv{subseteq_antisymmetric}.
\end{proof}

\begin{proposition}\label{dependent_product_coordinate}
    Let $f\in\dependentProduct{I}{A}$.
    Let $i\in I$.
    Then $f(i)\in A(i)$.
\end{proposition}
\begin{proof}
    Follows by \obv{dependent_product_iff}.
\end{proof}

\begin{proposition}\label{indexed_family_member_union}
    Suppose $A$ is a function on $I$.
    Let $i\in I$.
    Let $x\in A(i)$.
    Then $x\in\unions{\ran{A}}$.
\end{proposition}
\begin{proof}
    $A$ is a function and $I=\dom{A}$.
    Hence $\dom{A}=I$.
    Hence $A(i)\in\ran{A}$ by \auto{fun_ran_iff}.
    Follows by \auto{unions_iff}.
\end{proof}

\begin{proposition}\label{indexed_sum_iff}
    Suppose $A$ is a function on $I$.
    $(i,x)\in\indexedSum{I}{A}$ iff $i\in I$ and $x\in A(i)$.
\end{proposition}
\begin{proof}
    Show that if $(i,x)\in\indexedSum{I}{A}$, then
        $i\in I$ and $x\in A(i)$.
    \begin{subproof}
        Assume $(i,x)\in\indexedSum{I}{A}$.
        Hence $(i,x)\in I\times\unions{\ran{A}}$ and
            $\snd{(i,x)}\in A(\fst{(i,x)})$ by \obv{indexed_sum}.
        $\fst{(i,x)}=i$ and $\snd{(i,x)}=x$ by \obv{fst_eq,snd_eq}.
        Follows by \obv{times_tuple_elim}.
    \end{subproof}
    Show that if $i\in I$ and $x\in A(i)$, then
        $(i,x)\in\indexedSum{I}{A}$.
    \begin{subproof}
        Assume $i\in I$ and $x\in A(i)$.
        Hence $x\in\unions{\ran{A}}$
            by \obv{indexed_family_member_union}.
        Hence $(i,x)\in I\times\unions{\ran{A}}$
            by \obv{times_tuple_intro}.
        $\fst{(i,x)}=i$ and $\snd{(i,x)}=x$ by \obv{fst_eq,snd_eq}.
        Follows by \obv{indexed_sum}.
    \end{subproof}
    Follows by \auto.
\end{proof}

\begin{theorem}\label{dependent_product_inhabited_iff}
    Suppose $A$ is a function on $I$.
    $\dependentProduct{I}{A}$ is inhabited iff
        for every $i\in I$ we have $A(i)$ is inhabited.
\end{theorem}
\begin{proof}
    Show that if $\dependentProduct{I}{A}$ is inhabited, then
        for every $i\in I$ we have $A(i)$ is inhabited.
    \begin{subproof}
        Assume $\dependentProduct{I}{A}$ is inhabited.
        Take $f$ such that $f\in\dependentProduct{I}{A}$.
        For every $i\in I$ we have $f(i)\in A(i)$
            by \obv{dependent_product_coordinate}.
        Follows by \auto.
    \end{subproof}
    Show that if for every $i\in I$ we have $A(i)$ is inhabited, then
        $\dependentProduct{I}{A}$ is inhabited.
    \begin{subproof}
        Assume for every $i\in I$ we have $A(i)$ is inhabited.
        Let $f=\{(i,\choice{A(i)})\mid i\in I\}$.
        $f$ is a relation by \auto{relation}.
        $f$ is right-unique by \auto{rightunique,pair_eq_iff}.
        $\dom{f}=I$ by \auto{setext,dom_iff,pair_eq_iff}.
        For every $i\in I$ we have $\choice{A(i)}\in A(i)$
            by \obv{choice_in}.
        Hence for every $i\in I$ we have
            $\choice{A(i)}\in\unions{\ran{A}}$
            by \obv{indexed_family_member_union}.
        For every $i\in I$ we have $f(i)=\choice{A(i)}$
            by \obv{function_apply_intro}.
        Hence $f$ is a function to $\unions{\ran{A}}$ by \auto.
        Hence $f\in\funs{I}{\unions{\ran{A}}}$
            by \obv{funs_intro}.
        For every $i\in I$ we have $f(i)\in A(i)$ by \auto.
        Hence $f\in\dependentProduct{I}{A}$
            by \obv{dependent_product_iff}.
    \end{subproof}
    Follows by \auto.
\end{proof}

\begin{definition}\label{indexed_sum_injection}
    $\indexedSumInjection{I}{A}{i}
        =\{(x,(i,x))\mid x\in A(i)\}$.
\end{definition}

\begin{proposition}\label{indexed_sum_injection_type}
    Suppose $A$ is a function on $I$.
    Let $i\in I$.
    Then $\indexedSumInjection{I}{A}{i}
        \in\inj{A(i)}{\indexedSum{I}{A}}$.
\end{proposition}
\begin{proof}
    $\indexedSumInjection{I}{A}{i}$ is a relation by
        \auto{indexed_sum_injection,relation}.
    $\indexedSumInjection{I}{A}{i}$ is right-unique by
        \auto{indexed_sum_injection,rightunique,pair_eq_iff}.
    $\dom{\indexedSumInjection{I}{A}{i}}=A(i)$ by
        \auto{indexed_sum_injection,setext,dom_iff,pair_eq_iff}.
    For every $x\in A(i)$ we have $(i,x)\in\indexedSum{I}{A}$
        by \auto{indexed_sum_iff}.
    For every $x\in A(i)$ we have
        $(\indexedSumInjection{I}{A}{i})(x)=(i,x)$ by
        \obv{indexed_sum_injection,function_apply_intro}.
    Hence $\indexedSumInjection{I}{A}{i}
        \in\funs{A(i)}{\indexedSum{I}{A}}$ by \auto{funs_intro}.
    Show that $\indexedSumInjection{I}{A}{i}$ is injective.
    \begin{subproof}
        Show that for all $x,y$ such that
            $x\in\dom{\indexedSumInjection{I}{A}{i}}$ and
            $y\in\dom{\indexedSumInjection{I}{A}{i}}$ and
            $(\indexedSumInjection{I}{A}{i})(x)
                =(\indexedSumInjection{I}{A}{i})(y)$
            we have $x=y$.
        \begin{subproof}
            Fix $x,y$.
            Assume $x,y\in\dom{\indexedSumInjection{I}{A}{i}}$ and
                $(\indexedSumInjection{I}{A}{i})(x)
                    =(\indexedSumInjection{I}{A}{i})(y)$.
            Hence $x,y\in A(i)$.
            $(\indexedSumInjection{I}{A}{i})(x)=(i,x)$ and
                $(\indexedSumInjection{I}{A}{i})(y)=(i,y)$ by
                \obv{function_apply_intro}.
            Follows by \auto{pair_eq_iff}.
        \end{subproof}
        Follows by \auto{injective_function}.
    \end{subproof}
    Hence $\indexedSumInjection{I}{A}{i}
        \in\inj{A(i)}{\indexedSum{I}{A}}$
        by \auto{inj,injective_function}.
\end{proof}

\begin{proposition}\label{indexed_sum_injection_apply}
    Suppose $A$ is a function on $I$.
    Let $i\in I$ and $x\in A(i)$.
    Then $(\indexedSumInjection{I}{A}{i})(x)=(i,x)$.
\end{proposition}
\begin{proof}
    $\indexedSumInjection{I}{A}{i}
        \in\inj{A(i)}{\indexedSum{I}{A}}$ by
        \obv{indexed_sum_injection_type}.
    Hence $\indexedSumInjection{I}{A}{i}$ is a function by
        \auto{inj,funs_is_function}.
    $(x,(i,x))\in\indexedSumInjection{I}{A}{i}$ by
        \obv{indexed_sum_injection}.
    Follows by \obv{function_apply_intro}.
\end{proof}

\begin{proposition}\label{indexed_sum_injection_preimage_iff}
    $x\in\preimg{\indexedSumInjection{I}{A}{i}}{U}$ iff
        $x\in A(i)$ and $(i,x)\in U$.
\end{proposition}
\begin{proof}
    Follows by
        \auto{preimg_iff,indexed_sum_injection,pair_eq_iff}.
\end{proof}

\begin{proposition}\label{indexed_sum_injection_image}
    Suppose $U\subseteq A(i)$.
    Then $\img{\indexedSumInjection{I}{A}{i}}{U}=\{i\}\times U$.
\end{proposition}
\begin{proof}
    Show that every element of
        $\img{\indexedSumInjection{I}{A}{i}}{U}$ is an element of
        $\{i\}\times U$.
    \begin{subproof}
        Fix $p\in\img{\indexedSumInjection{I}{A}{i}}{U}$.
        Take $x$ such that $x\in U$ and
            $(x,p)\in\indexedSumInjection{I}{A}{i}$ by \obv{img_iff}.
        Hence $x\in A(i)$ and $p=(i,x)$ by
            \auto{indexed_sum_injection,pair_eq_iff}.
        $(i,x)\in\{i\}\times U$ by
            \auto{times_tuple_intro,singleton_iff}.
    \end{subproof}
    Hence $\img{\indexedSumInjection{I}{A}{i}}{U}
        \subseteq\{i\}\times U$ by \obv{subseteq}.

    Show that every element of $\{i\}\times U$ is an element of
        $\img{\indexedSumInjection{I}{A}{i}}{U}$.
    \begin{subproof}
        Fix $p\in\{i\}\times U$.
        Take $j,x$ such that $p=(j,x)$ by
            \auto{times_elem_is_tuple}.
        Hence $(j,x)\in\{i\}\times U$.
        Hence $j\in\{i\}$ and $x\in U$ by \obv{times_tuple_elim}.
        Hence $j=i$ by \obv{singleton_iff}.
        $x\in A(i)$ by \obv{subseteq}.
        $(x,(i,x))\in\indexedSumInjection{I}{A}{i}$ by
            \obv{indexed_sum_injection}.
        Hence $(i,x)\in\img{\indexedSumInjection{I}{A}{i}}{U}$ by
            \auto{img_iff}.
    \end{subproof}
    Hence $\{i\}\times U
        \subseteq\img{\indexedSumInjection{I}{A}{i}}{U}$ by
        \obv{subseteq}.
    Follows by \obv{subseteq_antisymmetric}.
\end{proof}

\begin{theorem}\label{indexed_sum_fibre_injection}
    Suppose $A$ is a function on $I$.
    Let $i\in I$.
    Then there exists $e\in\inj{A(i)}{\indexedSum{I}{A}}$.
\end{theorem}
\begin{proof}
    $\indexedSumInjection{I}{A}{i}
        \in\inj{A(i)}{\indexedSum{I}{A}}$ by
        \obv{indexed_sum_injection_type}.
\end{proof}

\begin{theorem}\label{indexed_union_injection_into_sum}
    Suppose $A$ is a function on $I$.
    Then there exists
        $e\in\inj{\unions{\ran{A}}}{\indexedSum{I}{A}}$.
\end{theorem}
\begin{proof}
    Let $U=\unions{\ran{A}}$.
    Let $R=\{p\in U\times I\mid \fst{p}\in A(\snd{p})\}$.
    $A$ is a function and $I=\dom{A}$.
    Hence $R$ is a relation by
        \auto{relation,times_elem_is_tuple}.
    Show that for every $x\in U$ there exists $i$ such that
        $x\mathrel{R}i$.
    \begin{subproof}
        Fix $x\in U$.
        Take $S\in\ran{A}$ such that $x\in S$ by \obv{unions_iff}.
        Take $i\in\dom{A}$ such that $A(i)=S$ by \obv{fun_ran_iff}.
        Hence $i\in I$.
        $(x,i)\in U\times I$ by \obv{times_tuple_intro}.
        $\fst{(x,i)}=x$ and $\snd{(x,i)}=i$ by \obv{fst_eq,snd_eq}.
        Hence $(x,i)\in R$.
    \end{subproof}
    Take $j$ such that $j$ is a function on $U$ and
        for every $x\in U$ we have $x\mathrel{R}j(x)$
        by \obv{axiom_of_choice}.
    Let $e=\{(x,(j(x),x))\mid x\in U\}$.
    $e$ is a relation by \auto{relation}.
    $e$ is right-unique by \auto{rightunique,pair_eq_iff}.
    $\dom{e}=U$ by \auto{setext,dom_iff,pair_eq_iff}.
    Show that for every $x\in U$ we have
        $(j(x),x)\in\indexedSum{I}{A}$.
    \begin{subproof}
        Fix $x\in U$.
        $(x,j(x))\in R$.
        Hence $(x,j(x))\in U\times I$ and
            $\fst{(x,j(x))}\in A(\snd{(x,j(x))})$.
        $\fst{(x,j(x))}=x$ and $\snd{(x,j(x))}=j(x)$
            by \obv{fst_eq,snd_eq}.
        Hence $j(x)\in I$ and $x\in A(j(x))$
            by \obv{times_tuple_elim}.
        Follows by \obv{indexed_sum_iff}.
    \end{subproof}
    For every $x\in U$ we have $e(x)=(j(x),x)$
        by \obv{function_apply_intro}.
    Hence $e\in\funs{U}{\indexedSum{I}{A}}$ by \auto{funs_intro}.
    Show that $e$ is injective.
    \begin{subproof}
        Show that for all $x,y$ such that $x\in\dom{e}$ and
            $y\in\dom{e}$ and $e(x)=e(y)$ we have $x=y$.
        \begin{subproof}
            Fix $x,y$.
            Assume $x\in\dom{e}$ and $y\in\dom{e}$ and $e(x)=e(y)$.
            Hence $x,y\in U$.
            $e(x)=(j(x),x)$ and $e(y)=(j(y),y)$
                by \obv{function_apply_intro}.
            Follows by \auto{pair_eq_iff}.
        \end{subproof}
        Follows by \auto{injective_function}.
    \end{subproof}
    Hence $e\in\inj{U}{\indexedSum{I}{A}}$
        by \auto{inj,injective_function}.
    Follows by \auto.
\end{proof}
